% eksperymentalne tłumaczenie na włoski

%

\section{Twierdzenie kosinusów} % lang-pl

\index{twierdzenie!kosinusów}% % lang-pl
\section{Teorema del coseno} % lang-it
\index{teorema!del coseno}% % lang-it

Twierdzenie kosinusów jest bardzo stare. % lang-pl

Euklides rozpatrzy je w Elementach osobno dla trójkątów rozwartokątnych (II.12) i ostrokątnych (II.13). % lang-pl
Il teorema del coseno è molto antico. % lang-it

Euclide lo tratta negli Elementi separatamente per i triangoli ottusangoli (II.12) e acutangoli (II.13). % lang-it
% TODO: https://en.wikipedia.org/wiki/Law_of_cosines The cases of obtuse triangles and acute triangles (corresponding to the two cases of negative or positive cosine) are treated separately, in Propositions II.12 and II.13:[1]
Perski matematyk Jamshid al-Kashi znajdzie wartość $2\pi$ z~dokładnością do szesnastu cyfr, będzie autorem najdokładniejszych tablic trygonometrycznych swoich czasów i poda równoważną postać wzoru, % lang-pl

Il matematico persiano Jamshid al-Kashi calcolerà il valore di $2\pi$ con precisione fino a sedici cifre, sarà autore delle tavole trigonometriche più precise del suo tempo e fornirà una forma equivalente della formula, % lang-it
\index[persons]{al-Kashi, Jamshid}
\begin{equation}
	c = \sqrt{(b - a \cos \gamma)^2 + (a \sin \gamma)^2}
\end{equation}
dla ostrego kąta $\gamma$. % lang-pl

Taka sama metoda rozwiązywania trójkątów pojawi się w Europie w 1464 roku, kiedy Regiomontanus opublikuje \emph{,,De triangulis omnimodis''} (czyli \emph{,,O trójkątach wszelkiego rodzaju''}). % lang-pl
per l'angolo acuto $\gamma$. % lang-it

Lo stesso metodo per risolvere i triangoli apparirà in Europa nel 1464, quando Regiomontano pubblicherà \emph{De triangulis omnimodis} (cioè \emph{Sui triangoli di ogni tipo}). % lang-it
\index[persons]{Regiomontanus}%
Współczesna forma twierdzenia to zasługa François Viète'a. % lang-pl

La forma moderna del teorema è merito di François Viète. % lang-it
\index[persons]{Viète, François}%

\begin{proposition}[twierdzenie kosinusów] % lang-pl
	W trójkącie o bokach długości $a, b, c$ z kątem $\gamma$ naprzeciwko krawędzi $c$ zachodzi % lang-pl
\begin{proposition}[teorema del coseno] % lang-it
	
	In un triangolo con lati di lunghezza $a, b, c$ e angolo $\gamma$ opposto al lato $c$ vale % lang-it
	\label{twierdzenie_kosinusow}%
	\begin{equation}
		c^2 = a^2 + b^2 - 2ab \cos \gamma.
	\end{equation}
	% https://en.wikipedia.org/wiki/Law_of_cosines
\end{proposition}

Znamy różne dowody tego twierdzenia: korzystające z twierdzenia Pitagorasa, trzech wysokości trójkąta, twierdzenia Ptolemeusza, geometrii koła albo twierdzenia sinusów. % lang-pl

We Francji twierdzenie kosinusów do późna będzie nazywane \emph{théorème d'Al-Kashi}, na cześć wspomnianego wyżej Persa. % lang-pl
Conosciamo varie dimostrazioni di questo teorema: basate sul teorema di Pitagora, sulle tre altezze del triangolo, sul teorema di Tolomeo, sulla geometria del cerchio o sul teorema dei seni. % lang-it
In Francia il teorema del coseno sarà a lungo chiamato \emph{théorème d'Al-Kashi}, in onore del suddetto persiano. % lang-it


Piszą o nim Guzicki \cite[s. 258]{guzicki_2021}, Hartshorne \cite[s. 185]{hartshorne2000} jako ćwiczenie. % lang-pl
Ne scrivono Guzicki \cite[s. 258]{guzicki_2021}, Hartshorne \cite[s. 185]{hartshorne2000} come esercizio. % lang-it

\begin{theorem}[Stewarta, 1746] % lang-pl
\begin{theorem}[di Stewart, 1746] % lang-it
\label{twierdzenie_stewarta}%
\index{twierdzenie!Stewarta}%
	W trójkącie $\triangle ABC$ o bokach długości $a, b, c$ poprowadzono czewianę z~wierzchołka $C$ do boku $AB$ o długości $d$, dzieląc ten bok na odcinki długości $n$ oraz $m$, jak na rysunku: % lang-pl
	
	In un triangolo $\triangle ABC$ con lati di lunghezza $a, b, c$ è stata tracciata la ceviana dal vertice $C$ al lato $AB$ di lunghezza $d$, dividendo tale lato in segmenti di lunghezza $n$ e $m$, come mostrato nella figura: % lang-it
	\begin{center}
\begin{comment}
    \begin{tikzpicture}[scale=.4]
        %\tkzInit[xmin=-0.5,xmax=6.5, ymin=-0.5,ymax=4.5]
        % \tkzClip
        \tkzDefPoint(0, 0){A}
		\tkzDefPoint(3.25, 3.25){d}

		\tkzDefPoint(6, 0){AB}
        \tkzDefPoint(10, 0){B}
        \tkzDefPoint(1, 7){C}
        \tkzDefPoint(35:4.75){CC}
		\tkzDrawSegments(C,AB)
        \tkzDrawPolygon[line width=0.3mm](A,B,C)

        \tkzLabelPoint[below left](A){$A$}
        \tkzLabelPoint[below right](B){$B$}
        \tkzLabelPoint[above](C){$C$}

        \tkzLabelPoint(d){$d$}

		\tkzDrawPoints[size=3,color=black,fill=black!80](A,B,C,AB)
		\tkzDrawSegment[dim={$\,\,c\,\,$,-16pt,transform shape}](A,B)
		\tkzDrawSegment[dim={$\,\,n\,\,$,-8pt,transform shape}](A,AB)
		\tkzDrawSegment[dim={$\,\,m\,\,$,-8pt,transform shape}](AB,B)
		\tkzDrawSegment[dim={$\,\,b\,\,$,8pt,transform shape,sloped}](A,C)
		\tkzDrawSegment[dim={$\,\,a\,\,$,-8pt,transform shape,sloped}](B,C)
    \end{tikzpicture}
\end{comment}
    \end{center}
	Wtedy % lang-pl
	
	Allora % lang-it
	\begin{equation}
		b^2 m + c^2 n = a (d^2 + mn).
	\end{equation}
\end{theorem}

Twierdzenie bez dowodu opublikuje Matthew Stewart w 1746 roku, chociaż Coxeter przypuści, że mogło być znane nawet Archimedesowi. % lang-pl

Il teorema, senza dimostrazione, sarà pubblicato da Matthew Stewart nel 1746, anche se Coxeter sospetta che fosse noto persino ad Archimede. % lang-it
\index[persons]{Stewart, Matthew}%
\index[persons]{Archimedes}% % lang-pl
\index[persons]{Archimede}% % lang-it
% Coxeter, H.S.M.; Greitzer, S.L. (1967), Geometry Revisited, New Mathematical Library #19, The Mathematical Association of America, ISBN 0-88385-619-0 strona 6
Dowody podadzą też Simpson, Euler, Carnot i pewnie jeszcze ktoś. % lang-pl

Dimostrazioni saranno fornite anche da Simpson, Euler, Carnot e probabilmente altri. % lang-it
\index[persons]{Simpson, Thomas}%
\index[persons]{Euler, Leonhard}%
\index[persons]{Carnot, L.M.N.}%
W przyszłości często pokazywać się będzie je jako zastosowanie twierdzenia kosinusów, tak jak Guzicki \cite[s. 265]{guzicki_2021}, Bogdańska, Neugebauer \cite[s. 90-91]{neugebauer_2018}; inne rozwiązanie skorzysta z~odcinków skierowanych. % lang-pl

In futuro saranno spesso presentati come applicazione del teorema del coseno, come Guzicki \cite[s. 265]{guzicki_2021}, Bogdańska, Neugebauer \cite[s. 90-91]{neugebauer_2018}; un'altra soluzione utilizzerà segmenti orientati. % lang-it
% TODO: Zetel, s. 31 podaje bez kosinusów, proste rachunki
Znalezienie go to ćwiczenie podane przez Evesa \cite[s. 58]{eves1_1972}. % lang-pl

Trovarlo è un esercizio proposto da Eves \cite[s. 58]{eves1_1972}. % lang-it

\begin{corollary}[twierdzenie Apoloniusza] % lang-pl
	W trójkącie $ABC$ poprowadzono środkową $AD$. % lang-pl
\begin{corollary}[teorema di Apollonio] % lang-it
	
	Wtedy suma pól kwadratów zbudowanych na bokach $AB$ i $AC$ jest dwa razy mniejsza od sumy pól kwadratów zbudowanych na środkowej $AD$ i połowie podstawy ($BD$): % lang-pl
	In un triangolo $ABC$ è stata tracciata la mediana $AD$. % lang-it
	
	Allora la somma delle aree dei quadrati costruiti sui lati $AB$ e $AC$ è la metà della somma delle aree dei quadrati costruiti sulla mediana $AD$ e sulla metà della base ($BD$): % lang-it
	\begin{equation}
		|AB|^2 + |AC|^2 = 2|BD|^2 + 2 |AD|^2.
	\end{equation}
	% https://en.wikipedia.org/wiki/Apollonius%27s_theorem
\end{corollary}

Twierdzenie przetrwa jako fakt VII.122 w zbiorach Pappusa z Aleksandrii, choć nie brak będzie przypuszczeń, że występowało wcześniej w zaginionym tekście \emph{,,Plane Loci''} Apoloniusza z Pergi. % lang-pl

(Tekst będzie wspomniany na stronie \pageref{plane_loci}). % lang-pl
Il teorema sopravvive come fatto VII.122 nelle collezioni di Pappo di Alessandria, anche se non mancheranno supposizioni che fosse presente in un testo perduto \emph{Plane Loci} di Apollonio di Perge. % lang-it

(Il testo sarà menzionato a pagina \pageref{plane_loci}). % lang-it

\begin{corollary}
	W trójkącie $\triangle ABC$ o bokach długości $a, b, c$ poprowadzono środkową oraz dwusieczną z~wierzchołka $C$. % lang-pl
	
	Długość środkowej wynosi % lang-pl
	In un triangolo $\triangle ABC$ con lati di lunghezza $a, b, c$ è stata tracciata la mediana e la bisettrice dall'angolo $C$. % lang-it
	
	La lunghezza della mediana è % lang-it
	\begin{equation}
		\sqrt{\frac{a^2 + b^2}{2} - \frac{c^2}{4}},
	\end{equation}
	zaś dwusieczna ma długość % lang-pl
	
	mentre la bisettrice ha lunghezza % lang-it
	\begin{equation}
		\frac{\sqrt{ab (a+b+c)(a+b-c)}}{a+b}.
	\end{equation}
\end{corollary}

%